\documentclass[12pt,letterpaper]{article}
\usepackage[utf8]{inputenc}

% Ajustamos los márgenes para corregir las advertencias de fancyhdr
\usepackage[letterpaper, margin=1in, bottom=3in, top=2.2in]{geometry}

% Cargamos primero el paquete fancyhdr
\usepackage{fancyhdr}
% Corregimos los parámetros headheight y footskip
\setlength{\headheight}{15pt} % Al menos 14.49998pt según la advertencia
\setlength{\footskip}{72pt}   % Al menos 71.60547pt según la advertencia

\usepackage{xcolor}
\usepackage{amsmath}
\usepackage{anyfontsize}
\usepackage{graphicx}
\usepackage{tikz} % Necesario para los elementos gráficos
\usepackage{listings} % Para resaltado de sintaxis
\usepackage{sourcesanspro}
\renewcommand{\familydefault}{\sfdefault}
\usepackage{fontawesome5}
\usepackage{titlesec}
\usepackage{setspace}
\usepackage{hyperref}
\usepackage{mdframed} % Para marcos más estables
\usepackage{qrcode} % Paquete para generar códigos QR

% Cargar tikz y la biblioteca de cajitas para mdframed
\usetikzlibrary{calc,shapes,positioning}

% IMPORTANT: Force the text color to be white for the entire document
\AtBeginDocument{\color{primaryColor}}

% Colors updated with better contrast
\definecolor{myblue}{cmyk}{1.0,0.8,0.05,0.20}
\definecolor{bgColor}{RGB}{15, 22, 36}  % Very dark blue, almost black background
\definecolor{primaryColor}{RGB}{255, 255, 255}  % White text - full white for better contrast
\definecolor{accentColor}{RGB}{255, 212, 59}  % Python Yellow
\definecolor{pythonBlue}{RGB}{120, 180, 255}  % Azul Python más claro para mejor contraste
\definecolor{secondaryColor}{RGB}{220, 230, 240}  % Lighter blue-gray for better contrast
\definecolor{terminalBg}{RGB}{22, 22, 22}  % Terminal background remains dark
\definecolor{terminalFrame}{RGB}{40, 40, 40}  % Terminal frame slightly darker
\definecolor{lineNumberColor}{RGB}{100, 100, 100}  % Color más sutil para los números de línea
\definecolor{dividerColor}{RGB}{60, 70, 90}  % Color más sutil para la línea divisoria

% Colores mejorados para el código con mayor contraste
\definecolor{codeTextColor}{RGB}{255, 255, 255}  % Blanco puro para el texto básico
\definecolor{codeKeywordColor}{RGB}{135, 206, 250}  % Azul cielo para palabras clave
\definecolor{codeCommentColor}{RGB}{144, 238, 144}  % Verde claro para comentarios
\definecolor{codeStringColor}{RGB}{255, 215, 135}  % Amarillo para strings
\definecolor{codeMethodColor}{RGB}{255, 160, 122}  % Naranja claro para métodos
\definecolor{codeFunctionColor}{RGB}{255, 215, 0}  % Dorado para funciones
\definecolor{codeNumberColor}{RGB}{255, 255, 150}  % Amarillo claro para números

% Set page background color
\pagecolor{bgColor}

% Hyperref configuration
\hypersetup{
    colorlinks=true,
    linkcolor=accentColor,
    filecolor=accentColor,
    urlcolor=accentColor,
}

% Configuración básica para listings que evita problemas con saltos de página
\lstset{
  language=Python,
  basicstyle=\normalsize\ttfamily\bfseries\color{codeTextColor}, % Tamaño reducido, en blanco y negrita
  backgroundcolor=\color{terminalBg},
  commentstyle=\color{codeCommentColor},
  keywordstyle=\color{codeKeywordColor},
  stringstyle=\color{codeStringColor},
  numberstyle=\color{lineNumberColor}, % Números de línea más sutiles
  breaklines=true,
  breakatwhitespace=true,
  tabsize=4,
  showstringspaces=false,
  frame=none,
  xleftmargin=15pt,
  xrightmargin=0pt,
  aboveskip=10pt,
  belowskip=10pt,
  numbers=left,
  numbersep=8pt,
  extendedchars=true,
  keepspaces=true,
  columns=flexible,
  lineskip=6pt, % Espacio entre líneas ligeramente reducido
  % Keywords de Python (palabras reservadas)
  morekeywords={self, yield, lambda, with, as, from, True, False, None, import, in, for, 
                if, elif, else, while, return, def, class, try, except, finally, raise, 
                break, continue, global, nonlocal, pass, assert, del, yield from},
  % Funciones incorporadas resaltadas de manera distinta
  emph={[2]print,range,sum,int,str,float,list,dict,set,tuple,next,len,type,map,filter,
         reduce,zip,enumerate,sorted,reversed,min,max,open,any,all},
  emphstyle={[2]\color{codeFunctionColor}}
}

% TERMINAL ESTILO MAC MEJORADA CON BOTONES A LA IZQUIERDA Y MENOS ESPACIO SUPERIOR
\newenvironment{macterminal}{%
    \begin{mdframed}[
        linecolor=terminalFrame,
        backgroundcolor=terminalBg,
        roundcorner=5pt,
        skipabove=10pt,
        skipbelow=10pt,
        linewidth=1pt,
        innertopmargin=10pt, % Reducido
        frametitle={%
            \tikz[baseline=(current bounding box.east), outer sep=0pt]{
                \fill[red!80!black] (0,0) circle (5pt);
                \fill[yellow!80!black] (0.7,0) circle (5pt);
                \fill[green!70!black] (1.4,0) circle (5pt);
            }
        },
        frametitlealignment=\raggedright, % Alineado a la izquierda
        frametitleaboveskip=8pt, % Espacio entre el título y el borde superior
        frametitlebelowskip=0pt, % Reducir el espacio entre el título y el contenido
    ]
    % No necesitamos más espacio vertical aquí
}{%
    \end{mdframed}%
}

% Utilities
\newcommand{\verspace}{\vspace{10pt}}

% Section styling - Esquema más elegante y mejorado contraste
\titleformat{\section}
  {\LARGE\bfseries\color{primaryColor}} % Secciones principales en blanco para mejor contraste y más grandes
  {\thesection. }
  {0pt}
  {}
  []

\titleformat{\subsection}
  {\Large\bfseries\color{accentColor}} % Subsecciones en amarillo Python y más grandes
  {\thesubsection. }
  {0pt}
  {}
  []

\titleformat{\subsubsection}
  {\large\bfseries\color{pythonBlue}} % Subsubsecciones en azul Python más claro y más grandes
  {\thesubsubsection. }
  {0pt}
  {}
  []

% Spacing for sections
\titlespacing*{\section}{0pt}{20pt}{10pt}
\titlespacing*{\subsection}{0pt}{15pt}{7pt}
\titlespacing*{\subsubsection}{0pt}{10pt}{5pt}

% Configure header and footer
\pagestyle{fancy}
\fancyhf{}

% Crear un encabezado elegante para todas las páginas (excepto la portada)
\fancyhead[L]{
    \begin{tikzpicture}[remember picture, overlay]
        % Rectángulo de fondo para la sección Python
        \fill[pythonBlue, rounded corners=3pt] (0,0) rectangle (2.2cm,0.7cm);
        % Texto Python
        \node[text=primaryColor, font=\bfseries] at (1.1cm,0.35cm) {PYTHON};
        
        % Pequeño separador
        \fill[secondaryColor] (2.35cm,0.1cm) rectangle (2.40cm,0.6cm);
        
        % Texto Concurrency & Parallelism con el color de acento
        \node[text=accentColor, font=\bfseries, anchor=west] at (2.35cm,0.35cm) {CONCURRENCY (PART 1)};
    \end{tikzpicture}
}

% Línea de encabezado
\renewcommand{\headrulewidth}{0pt}


% Configure footer with profile info
\fancyfoot[C]{
    \vspace*{0.2cm}
    \noindent
    \begin{minipage}{\textwidth}
        \begin{flushleft}
            % Profile image placeholder with TikZ
            \raisebox{0.7cm}{
            \begin{tikzpicture}[baseline]
                \path[fill=bgColor] (0,0) circle (0.8cm);
                \clip (0,0) circle (0.8cm);
                \node at (0,0) {
                    % Placeholder for profile image
                    \includegraphics[width=1.6cm,height=1.6cm]{profile-image.jpeg}
                };
            \end{tikzpicture}
            }
            \begin{minipage}[b]{0.8\textwidth}
                % Name
                {\large\bfseries\color{primaryColor}Alejandro Sánchez Yalí}
                
                % Professional description
                \par\vspace{1pt}
                {\small\color{secondaryColor}Software Developer | AI \& Blockchain Enthusiast}
                
                % Contact
                \par\vspace{1pt}
                {\small\color{accentColor}\faGlobe\hspace{5pt}\color{secondaryColor}www.asanchezyali.com}
            \end{minipage}
        \end{flushleft}
        \vspace{8pt} % Space at bottom
    \end{minipage}
}

% Estilo para la portada (sin encabezado ni pie de página)
\fancypagestyle{plain}{
    \fancyhf{}
    \renewcommand{\headrulewidth}{0pt}
    \renewcommand{\footrulewidth}{0pt}
}

% Bullet styling for lists
\renewcommand{\labelitemi}{\textcolor{accentColor}{$\bullet$}} % Primer nivel de viñetas en amarillo (mejor contraste)
\renewcommand{\labelitemii}{\textcolor{pythonBlue}{$\circ$}} % Segundo nivel de viñetas en azul

% Aumentar globalmente el tamaño de texto un 20%
\usepackage{relsize}
\AtBeginDocument{\relsize{1}} % Incrementa la fuente en aproximadamente un 20%

% Comando simplificado para la etiqueta Python
\newcommand{\languagetag}[1]{
    \begin{tikzpicture}[baseline]
        \node[fill=pythonBlue, text=primaryColor, rounded corners=5pt, inner sep=7pt] {
            {\normalsize\textbf{#1}}
        };
    \end{tikzpicture}
}

% Alternativa usando una versión modificada de la URL
\newcommand{\elegantqr}[2]{
    \qrcode[height=2.5cm]{#1}
    \\[0.1cm]
    {\hspace{0.2cm}\color{primaryColor}\small #2\par}
}

% Definir el título personalizado con alineación exacta para todos los elementos
\newcommand{\titlepagecontents}{%
    \vspace*{2cm}
    % Asegurar que todos los elementos estén alineados a la izquierda
    \begin{flushleft}
    \languagetag{Python}\\[0.4cm]
    {\fontsize{48}{52}\bfseries\color{primaryColor}Python \color{accentColor}Concurrency\par}
    \vspace{0.3cm}
    {\fontsize{18}{52}\color{secondaryColor}Part 1: Understanding the GIL, Threading and Multiprocessing in Python\par}
    \vspace{0.3cm}
    {\color{secondaryColor}\today\par}
    \vspace{2cm}
    % El QR se alineará a la izquierda aunque internamente esté centrado
    \elegantqr{https://github.com/asanchezyali/social-media-posts/tree/main/Python/Concurrency}{Source Code}
    \end{flushleft}
}


% En la página del título, aplicar el contenido sin márgenes adicionales
\makeatletter
\renewcommand{\maketitle}{%
    \begin{titlepage}
    \setlength{\parindent}{0pt}
    \setlength{\leftskip}{0pt}
    \titlepagecontents
    \end{titlepage}
}
\makeatother

\newcommand{\finalpagecontents}{%
    \vspace*{4cm}
    \begin{center}
        \begin{tikzpicture}
            % Title text
            \node[text width=14cm, align=center] at (0,0) {
                {\fontsize{48}{52}\bfseries\color{primaryColor}Ready to \color{accentColor}Level Up\\\color{primaryColor}Your Python?\par}
            };
            
            % Add vertical space
            \node at (0,-3) {};
            
            % Question text
            \node[text width=14cm, align=center] at (0,-5) {
                {\fontsize{22}{26}\color{secondaryColor}You've mastered powerful concurrency techniques.\\\vspace{0.5cm}Apply these patterns in your next project\\\vspace{0.5cm}and watch your code become faster and more efficient!\par \vspace{1cm}
                \color{accentColor}Stay tuned for Part 2...\\\vspace{0.4cm}}
                
            };
        \end{tikzpicture}
    \end{center}
}

\begin{document}
\color{primaryColor} % Explicitly set text color to white for the entire document

\begin{titlepage}
    % Eliminamos todos los espacios adicionales para asegurar la alineación
    \titlepagecontents
\end{titlepage}

\section{Understanding Concurrency and Parallelism}

Concurrency and parallelism are fundamental concepts in Python for improving application performance, but they serve different purposes and work in distinct ways:

\begin{itemize}
    \item \textbf{\textcolor{pythonBlue}{Concurrency:}} Dealing with multiple tasks by switching between them (not necessarily simultaneously)
    \item \textbf{\textcolor{pythonBlue}{Parallelism:}} Executing multiple tasks truly simultaneously using multiple processors
    \item \textbf{\textcolor{pythonBlue}{Performance Impact:}} Proper implementation can dramatically speed up I/O-bound or CPU-bound operations
\end{itemize}

\subsection{When to Use Each Approach}

Choosing the right approach depends on your task type:

\begin{itemize}
    \item \textbf{\textcolor{pythonBlue}{For I/O-Bound Tasks}} (waiting for external resources):
    \begin{itemize}
        \item Network requests, file operations, database queries
        \item \textbf{Best approach:} Threading or Asyncio
    \end{itemize}
    \item \textbf{\textcolor{pythonBlue}{For CPU-Bound Tasks}} (heavy computations):
    \begin{itemize}
        \item Data processing, image manipulation, numerical calculations
        \item \textbf{Best approach:} Multiprocessing
    \end{itemize}
\end{itemize}

\section{The Global Interpreter Lock (GIL)}

\subsection{Understanding the GIL's Impact}

The Global Interpreter Lock (GIL) is a mutex (lock) that protects access to Python objects, preventing multiple threads from executing Python bytecode simultaneously in a single process.

\begin{itemize}
    \item \textbf{\textcolor{pythonBlue}{Single-Threaded Execution:}} Only one thread can execute Python code at any given time
    \item \textbf{\textcolor{pythonBlue}{I/O Release:}} The GIL is released during I/O operations, allowing other threads to run
    \item \textbf{\textcolor{pythonBlue}{CPU Limitation:}} CPU-bound threads still time-share a single CPU core
    \item \textbf{\textcolor{pythonBlue}{Multiprocessing Bypass:}} Multiprocessing avoids the GIL by using multiple processes
\end{itemize}

Think of the GIL as a talking stick in a meeting—only the person holding it can speak (execute code), and it must be passed around even when multiple people (threads) want to talk simultaneously.

\subsection{GIL Workarounds}

\begin{itemize}
    \item \textbf{\textcolor{pythonBlue}{Use multiprocessing}} for CPU-intensive tasks
    \item \textbf{\textcolor{pythonBlue}{Leverage asyncio}} for I/O-bound operations
    \item \textbf{\textcolor{pythonBlue}{Offload to C extensions}} that release the GIL
    \item \textbf{\textcolor{pythonBlue}{Consider alternative implementations}} like PyPy for specific use cases
\end{itemize}

\section{Threading in Python}

\subsection{Thread Basics}

Threads share memory space but execute concurrently:

\begin{macterminal}
\begin{lstlisting}
import threading
import time

def worker(name, delay):
    """A simple worker function that demonstrates a thread"""
    print(f"Worker {name}: Starting")
    for i in range(3):
        time.sleep(delay)
        print(f"Worker {name}: Step {i+1}")
    print(f"Worker {name}: Finished")

# Create thread objects
thread1 = threading.Thread(target=worker, args=("A", 0.5))
thread2 = threading.Thread(target=worker, args=("B", 0.8))

# Start the threads
print("Main thread: Starting workers")
thread1.start()
thread2.start()

# Continue with other work in the main thread
print("Main thread: Doing other work")
time.sleep(1)
print("Main thread: Work complete")

# Wait for worker threads to finish
thread1.join()
thread2.join()
print("Main thread: All workers finished")
\end{lstlisting}
\end{macterminal}

\subsection{Thread Safety and Locks}

When multiple threads access shared data, synchronization is essential:

\begin{macterminal}
\begin{lstlisting}
import threading
import time

class Counter:
    def __init__(self):
        self.value = 0
        self.lock = threading.Lock()
        
    def increment(self):
        with self.lock:  # Thread-safe access to shared data
            current = self.value
            time.sleep(0.001)  # Simulate work
            self.value = current + 1
    
    def get_value(self):
        with self.lock:
            return self.value

def increment_counter(counter, count):
    for _ in range(count):
        counter.increment()

# Create a shared counter
counter = Counter()

# Create threads to increment the counter
threads = []
num_threads = 10
increments_per_thread = 100

for _ in range(num_threads):
    t = threading.Thread(
        target=increment_counter, 
        args=(counter, increments_per_thread)
    )
    threads.append(t)
    t.start()

# Wait for all threads to complete
for t in threads:
    t.join()

# Check the result
expected = num_threads * increments_per_thread
actual = counter.get_value()
print(f"Expected count: {expected}")
print(f"Actual count: {actual}")
print(f"Thread-safe: {expected == actual}")
\end{lstlisting}
\end{macterminal}

\section{Multiprocessing for Parallel Computing}

\subsection{Process Basics}

Processes have separate memory spaces and bypass the GIL:

\begin{macterminal}
\begin{lstlisting}
import multiprocessing as mp
import time
import os

def cpu_bound_task(number):
    """A CPU-intensive function that benefits from parallelization"""
    print(f"Process {os.getpid()} processing {number}")
    # Simulate CPU-intensive calculation
    result = 0
    for i in range(10**7):  # 10 million iterations
        result += i * number
    return result

if __name__ == "__main__":
    numbers = [1, 2, 3, 4, 5, 6, 7, 8]
    
    # Sequential processing
    start_time = time.time()
    sequential_results = [cpu_bound_task(n) for n in numbers]
    sequential_time = time.time() - start_time
    print(f"Sequential processing took {sequential_time:.2f} seconds")
    
    # Parallel processing
    start_time = time.time()
    with mp.Pool(processes=mp.cpu_count()) as pool:
        parallel_results = pool.map(cpu_bound_task, numbers)
    parallel_time = time.time() - start_time
    print(f"Parallel processing took {parallel_time:.2f} seconds")
    print(f"Speedup: {sequential_time / parallel_time:.2f}x")
\end{lstlisting}
\end{macterminal}

\subsection{Sharing Data Between Processes}

Unlike threads, processes require special mechanisms to share data:

\begin{macterminal}
\begin{lstlisting}
import multiprocessing as mp

def update_shared_dict(shared_dict, key, value):
    """Update a value in a shared dictionary"""
    shared_dict[key] = value
    print(f"Process {mp.current_process().name} updated shared_dict[{key}] = {value}")

def sum_shared_array(shared_array, result_queue):
    """Calculate sum of a shared array and put result in a queue"""
    total = sum(shared_array)
    result_queue.put(total)
    print(f"Process {mp.current_process().name} calculated sum: {total}")

if __name__ == "__main__":
    # Create a manager to coordinate shared objects
    with mp.Manager() as manager:
        # Create shared objects
        shared_dict = manager.dict()  # Shared dictionary
        shared_array = manager.list(range(10))  # Shared list
        result_queue = manager.Queue()  # Shared queue
        
        # Create processes
        processes = []
        
        # Processes to update the shared dictionary
        for i in range(5):
            p = mp.Process(
                target=update_shared_dict,
                args=(shared_dict, f"key_{i}", i*10)
            )
            processes.append(p)
        
        # Process to calculate sum of shared array
        p = mp.Process(
            target=sum_shared_array,
            args=(shared_array, result_queue)
        )
        processes.append(p)
        
        # Start all processes
        for p in processes:
            p.start()
            
        # Wait for all processes to finish
        for p in processes:
            p.join()
        
        # Retrieve and display results
        print("\nShared dictionary:", dict(shared_dict))
        print("Sum from queue:", result_queue.get())
\end{lstlisting}
\end{macterminal}

\section{Synchronization Tools}

\subsection{Thread Synchronization Primitives}

Python provides several mechanisms for coordinating threads:

\begin{itemize}
    \item \textbf{\textcolor{pythonBlue}{Lock:}} The most basic synchronization primitive, prevents simultaneous access
    \item \textbf{\textcolor{pythonBlue}{RLock:}} Reentrant lock that can be acquired multiple times by the same thread
    \item \textbf{\textcolor{pythonBlue}{Semaphore:}} Controls access to a shared resource by multiple threads
    \item \textbf{\textcolor{pythonBlue}{Event:}} Signals between threads when a condition is met
    \item \textbf{\textcolor{pythonBlue}{Condition:}} More sophisticated signaling mechanism for complex coordination
    \item \textbf{\textcolor{pythonBlue}{Barrier:}} Makes threads wait until a specific number of them reach a common point
\end{itemize}

\begin{macterminal}
\begin{lstlisting}
import threading
import time

# Example: Using a semaphore to limit concurrent resource access
class ConnectionPool:
    def __init__(self, max_connections=3):
        self.semaphore = threading.Semaphore(max_connections)
        
    def get_connection(self):
        self.semaphore.acquire()
        return self  # Return a connection object
        
    def release_connection(self):
        self.semaphore.release()
        
    def __enter__(self):
        return self.get_connection()
        
    def __exit__(self, exc_type, exc_val, exc_tb):
        self.release_connection()

# Example usage
def worker(pool, worker_id):
    with pool:
        print(f"Worker {worker_id}: Acquired connection")
        time.sleep(1)  # Simulate work
        print(f"Worker {worker_id}: Released connection")

# Create a connection pool with max 3 connections
pool = ConnectionPool(3)

# Create and start 10 worker threads
threads = []
for i in range(10):
    t = threading.Thread(target=worker, args=(pool, i))
    threads.append(t)
    t.start()

# Wait for all threads to complete
for t in threads:
    t.join()
\end{lstlisting}
\end{macterminal}

\section{Conclusion (Part 1)}

In this first part of our exploration of concurrency and parallelism in Python, we have covered the essential fundamentals:

\begin{itemize}
    \item We understood the difference between concurrency (handling multiple tasks by alternating between them) and parallelism (executing multiple tasks simultaneously).
    
    \item We analyzed the Global Interpreter Lock (GIL) and its implications for concurrent programming in Python.
    
    \item We explored the \texttt{threading} module for concurrent programming oriented to I/O operations.
    
    \item We implemented proper synchronization to protect data shared between threads.
    
    \item We studied the \texttt{multiprocessing} module for true parallel execution.
    
    \item We learned about communication mechanisms between processes.
\end{itemize}

These tools are fundamental for developing efficient Python applications that handle multiple tasks. In the second part, we will dive deeper into \texttt{asyncio}, advanced patterns, and strategies for combining different concurrency models.

\section{References (Part 1)}
\begin{itemize}
    \item Python Documentation. \textit{threading - Thread-based parallelism}. \href{https://docs.python.org/3/library/threading.html}{Link}
    
    \item Python Documentation. \textit{multiprocessing - Process-based parallelism}. \href{https://docs.python.org/3/library/multiprocessing.html}{Link}
    
    \item Real Python. \textit{Python's GIL: A Guide to How It Impacts Performance}. \href{https://realpython.com/python-gil/}{Link}
    
    \item Python Discord. \textit{Concurrency, parallelism, threading and multiprocessing}. \href{https://pythondiscord.com/pages/resources/guides/concurrency-parallelism/}{Link}

    \item David Beazley. \textit{Understanding the Python GIL}. \href{https://www.dabeaz.com/python/UnderstandingGIL.pdf}{Link}
    
    \item Brett Slatkin. \textit{Effective Python}, Item 52-56: Concurrency and Parallelism. Addison-Wesley.
    
    \item Translated, Edited and written in collaboration with AI.
\end{itemize}

\section{Explore My Other Posts}
\vspace{10pt}
\begin{tikzpicture}
  % Section background - increased height slightly to ensure all content fits
  \fill[color=terminalBg, rounded corners=5pt] (0,0) rectangle (\textwidth,-5);
  
  % Calculate width available for text (total width minus QR code width minus margins)
  \pgfmathsetmacro{\availableWidth}{\textwidth-4cm}
  
  % Left column content
  \begin{scope}
    % Section title
    \node[text=accentColor, font=\Large\bfseries, anchor=north west] at (0.5cm,-0.5cm) {Enjoyed This Content?};
    
    % Subtitle
    \node[text=primaryColor, font=\normalsize, anchor=north west] at (0.5cm,-1.2cm) {Don't miss my previous post about:};
    
    % Python post title - with highlighted style and limited width
    \node[text=secondaryColor, font=\large\bfseries, anchor=north west, text width=\availableWidth] at (0.5cm,-1.9cm) {\color{accentColor}Python Context Managers: \color{primaryColor}Elegant Resource Management with the «{\color{accentColor}with}» Statement};
    
    % Brief description of the post - with limited width and adjusted position
    \node[text=secondaryColor, font=\normalsize, text width=\availableWidth, anchor=north west] at (0.5cm,-3.1cm) {
      Discover how Python Context Managers can help you elegantly manage resources, 
      prevent memory leaks, and write cleaner code.
    };
  \end{scope}
  
  % QR Code on the right - vertically centered
  \node[anchor=center] at ({\textwidth-1.5cm}, {-2.5cm}) {
    \qrcode[height=2.2cm]{https://www.linkedin.com/feed/update/urn:li:activity:7307438566157021186/}
  };
  
\end{tikzpicture}
\vspace{5pt}

\clearpage
\thispagestyle{empty}
\finalpagecontents

\end{document}